% Generated draft: requires mathematical and source review.
\documentclass[11pt]{article}
\usepackage[margin=25mm]{geometry}
\usepackage{fontspec}
\setmainfont{Latin Modern Roman}
\usepackage{amsmath,amssymb,mathtools,unicode-math}
\setmathfont{Latin Modern Math}
\usepackage{enumitem,needspace,xurl,hyperref}
\setlength{\parindent}{0pt}
\setlength{\parskip}{5pt}
\setlength{\emergencystretch}{3em}
\newcommand{\N}{\mathbb{N}}
\newcommand{\Z}{\mathbb{Z}}
\newcommand{\Q}{\mathbb{Q}}
\newcommand{\R}{\mathbb{R}}
\newcommand{\C}{\mathbb{C}}
\newcommand{\F}{\mathbb{F}}
\newcommand{\A}{\mathbb{A}}
\newcommand{\PP}{\mathbb{P}}
\newcommand{\E}{\mathbb{E}}
\newcommand{\eps}{\varepsilon}
\newcommand{\dd}{\,\mathrm{d}}
\begin{document}
\section*{UnsolvedMath ID: 30004600}
\textbf{Source problem code:} catalogue entry
\subsection{Definitions and mathematical statement}
Let \(\mathscr S\) denote the intended class of simplicial hyperplane arrangements. In the usual real setting, an arrangement is a collection of hyperplanes, its chambers are the connected components of their complement, and simpliciality means that its chambers are simplicial cones in the central essential case. The supplied record does not specify finiteness, centrality, essentiality, ambient dimension, or the treatment of nonessential or affine arrangements; this description therefore does not fix the scope of \(\mathscr S\).

Write \(\operatorname{CN}(\mathcal A)\), \(\operatorname{Free}(\mathcal A)\), \(\operatorname{Tight}(\mathcal A)\), and \(\operatorname{Sh}(\mathcal A)\) for congruence normality, freeness, tightness, and the shard digraph, respectively, in the conventions intended by the cited report. These are source-dependent terms: their definitions are not supplied in the record. In particular, the record does not specify any choice of base chamber or associated lattice used to define congruence normality and the shard digraph. Thus the following is a statement of the requested research targets, with unresolved conventions retained explicitly.

\begin{enumerate}
\item Classify the arrangements \(\mathcal A\in\mathscr S\) for which \(\operatorname{CN}(\mathcal A)\) holds. The equivalence relation intended for this classification is not specified.
\item Determine the relationship between \(\operatorname{CN}(\mathcal A)\) and \(\operatorname{Free}(\mathcal A)\) for \(\mathcal A\in\mathscr S\), including which implications hold and under what hypotheses. No implication is asserted by the supplied statement.
\item Identify useful properties of \(\operatorname{Sh}(\mathcal A)\) that can be described or interpreted in terms of the arrangement \(\mathcal A\). The requested properties and the criterion of usefulness are not specified.
\item Determine the smallest number of hyperplanes in a tight arrangement that is not congruence normal. More explicitly, if \(\mathscr T\) denotes the intended class of finite tight arrangements, determine
\[
\min\bigl\{\lvert\mathcal A\rvert:\mathcal A\in\mathscr T,\ \neg\operatorname{CN}(\mathcal A)\bigr\},
\]
provided this set is nonempty, and determine whether such arrangements exist. Here \(\lvert\mathcal A\rvert\) is the number of hyperplanes. The record does not specify whether \(\mathscr T\) is restricted to simplicial arrangements, which ranks are allowed, or the precise meaning of tightness.
\item Determine the behavior of congruence normality under restrictions and other arrangement operations. For a linear arrangement and a subspace \(X\), the usual restriction is
\[
\mathcal A^{X}=\{H\cap X:H\in\mathcal A,\ X\not\subseteq H\},
\]
with repeated hyperplanes identified. The record does not specify the allowed subspaces \(X\), whether this convention is intended, the other operations to consider, or hypotheses ensuring that the resulting arrangement lies in the domain of congruence normality. Preservation, reflection, and necessary additional hypotheses remain questions rather than asserted properties.
\end{enumerate}

A complete formal statement requires the missing scope and source-specific definitions. The supplied label \emph{open} is a curation label, not an independently verified mathematical conclusion.
\subsection{Short English statement}
Classify congruence-normal simplicial arrangements, relate congruence normality to freeness, interpret properties of the shard digraph through arrangement theory, find the smallest tight counterexample to congruence normality, and study restrictions and other operations. The source leaves several definitions and scope choices unspecified, so this statement-only draft requires human mathematical review.
\subsection{Sources}
\begin{itemize}
\item Local UnsolvedMath registry: \texttt{data/unsolvedmath/problems.json} and supplementary registry \texttt{catalog/problems.json}; supplied complete record with permanent UnsolvedMath ID \(30004600\), title \emph{Congruence Normality in Simplicial Arrangements}, and exact source problem code \texttt{catalogue entry}. The record belongs to the Oberwolfach Reports Open Problems collection and supplies proposed year 2021; no individual proposer is named.
\item Supplied report citation: \emph{Logarithmic Vector Fields and Freeness of Divisors and Arrangements: New perspectives and applications} (2022), Oberwolfach Reports, pp.\ 225--300, DOI \texttt{10.4171/owr/2021/5}. \url{https://doi.org/10.4171/owr/2021/5}. The background describes these reports as documenting workshops and their open-problem discussions. This authoritative report reference is supplied; its contents have not been independently consulted here.
\item Supplied literature source: Cuntz--Elia--Labbe, \emph{Beitr\"age zur Algebra und Geometrie}. \url{https://doi.org/10.1007/s00026-021-00555-2}. The record supplies no article title or more detailed bibliographic information.
\item Supplied literature source: \emph{Oberwolfach Report 5/2021}. \url{https://doi.org/10.4171/owr/2021/5}.
\item Literature assessment dated exactly \texttt{2026-08-22}: the supplied curation labels the problem open and states that known finite lists and rank-three classifications do not settle completeness of congruence-normal simplicial arrangements or the listed restriction and tightness questions. This is supplied dated curation derived from the OWR material, not independently verified literature assessment. The background expressly requests independent verification of time-sensitive status and cited claims.
\item The supplied record labels its statement \texttt{corrected\_verified} and reports a manual semantic, typographic, and LaTeX review claiming self-containment and preservation of the target. That is a supplied review claim; it does not resolve the formulation gaps identified in this draft.
\end{itemize}
\end{document}
